% ============================================================
%  EPIGRAM — Week 1, Topic 1: Foundations of Probability Modeling
%  Requires: epigram.sty (place in the same folder)
%  Compile:  pdflatex (twice for TOC)
% ============================================================
\documentclass[11pt]{article}
\usepackage{epigram}
\epigramsetup{1}{1}{Foundations of Probability Modeling}

% ── ID toggle ────────────────────────────────────────────────
\showidsfalse   % <── student PDF  (comment out to switch)
%\showidstrue   % <── backend view (uncomment to switch)

% ============================================================
\begin{document}

\epigramtitle
  {Counting $\cdot$ Inclusion-Exclusion $\cdot$ Geometric Probability $\cdot$ Combinatorics}
  {6}{7}
  {90--120 minutes}
  {None}


% ── Learning Objectives ──────────────────────────────────────
\begin{tcolorbox}[
  breakable,
  colback=EpiNavy!4,
  colframe=EpiNavy,
  arc=4pt,
  left=12pt, right=12pt, top=10pt, bottom=10pt,
  toptitle=5pt, bottomtitle=5pt,
  colbacktitle=EpiNavy,
  fonttitle=\bfseries\normalsize\color{white},
  title={{\large\strut} By the end of this topic, you will be able to\ldots},
]
\begin{enumerate}[leftmargin=16pt, itemsep=4pt, topsep=2pt]
  \item \textbf{Set up any probability problem} from scratch --- define the
        sample space $\Omega$, identify the event $A$, and compute
        $\Prob(A)$ for uniform probability spaces.
  \item \textbf{Apply complementary counting} to handle ``at least one''
        problems, including the Birthday Paradox and its generalisation.
  \item \textbf{Use the Inclusion-Exclusion Principle} to compute
        probabilities involving overlapping events and consecutive patterns.
  \item \textbf{Count arrangements under constraints} --- circular
        permutations, non-adjacent selections, and the block trick for
        consecutive groups.
  \item \textbf{Derive and apply the derangement formula}, and know that
        $\Prob(\text{derangement}) \approx 1/e$ for $n \geq 5$.
  \item \textbf{Solve geometric probability problems} by mapping events
        onto regions and computing areas --- including chord intersection
        and semicircle containment.
  \item \textbf{Recognise and exploit symmetry} as the fastest path to an
        answer in interview settings.
\end{enumerate}
\vspace{4pt}
{\small\color{EpiDarkGray}\textit{%
  Problems 1--6 (\freelabel) build all of the above skills.
  Problems 7--11 and A1--A2 (\premiumlabel) are the level you need to
  crack final rounds at Citadel, Jane Street, and Two Sigma.}}
\end{tcolorbox}
\vspace{8pt}

\section{Introduction: What Is Probability?}

Probability is the mathematics of uncertainty. In quantitative finance and
trading, every model, every bet, and every risk estimate rests on
probabilistic foundations. The goal of this topic is to build the
\textit{problem-solving toolkit} that lets you crack interview questions
under pressure in 5--10 minutes.

\begin{conceptbox}{The Three Ingredients of Any Probability Problem}
\begin{enumerate}
  \item \textbf{Sample space} $\Omega$: the set of all possible outcomes.
  \item \textbf{Event} $A \subseteq \Omega$: outcomes you care about.
  \item \textbf{Probability measure} $\Prob$: $\Prob(\Omega)=1$, $\Prob(\emptyset)=0$, countable additivity.
\end{enumerate}
For \textbf{uniform} spaces:
$\Prob(A) = |A|/|\Omega| = \text{favourable}/\text{total}$.
\end{conceptbox}

\textbf{Interview strategy:} (1) Identify $\Omega$. (2) Count $|A|$ — directly or via
complementary counting / IEP. (3) Sanity-check boundary cases.

\section{Counting Fundamentals}

\subsection{Permutations and Combinations}

\begin{conceptbox}{Key Formulas}
\begin{align*}
  \text{Permutations (ordered):}\quad
    & P(n,k) = \frac{n!}{(n-k)!} \\[4pt]
  \text{Combinations (unordered):}\quad
    & \dbinom{n}{k} = \frac{n!}{k!\,(n-k)!}
\end{align*}
\textit{Order matters $\Rightarrow$ permutation. Order doesn't $\Rightarrow$ combination.}
\end{conceptbox}

\subsection{Warm-Up}

\begin{workedbox}{1}{\freelabel\quad\difficultytitle{easy}}
\textbf{Majority Heads in Three Coin Tosses.}
Flip a fair coin three times. What is the probability of getting more heads than tails?
\end{workedbox}

\begin{solutionbox}
$|\Omega|=8$. Favourable (\#H $\geq 2$): $\{$HHH, HHT, HTH, THH$\}$ = 4 outcomes.
\[
  \Prob = \frac{4}{8} = \frac{1}{2}
\]
\textbf{Slick:} Odd tosses $\Rightarrow$ no ties $\Rightarrow$ by symmetry
$\Prob(\#H>\#T)=\Prob(\#T>\#H)=\tfrac{1}{2}$.

\textbf{Key takeaway:} For any odd number of fair coin tosses, majority heads has
probability $\tfrac{1}{2}$ by symmetry — no calculation needed.
\end{solutionbox}

\begin{freeproblem}{2}{\freelabel}{\difficultytitle{medium}}{quant\_interview\_fundamental\_qdsi\_p40}
\textbf{Parity of Heads in Coin Tosses.}
Flip 10 fair coins. What is the probability the number of heads is even?

\textit{Hint: try the algebraic trick $(p+(1-p))^n \pm ((1-p)-p)^n$.}
\end{freeproblem}

\begin{solutionbox}
\textbf{Method 1 — Direct:}
$|A|=\dbinom{10}{0}+\dbinom{10}{2}+\cdots+\dbinom{10}{10}=512$,
$\Prob=512/1024=\tfrac{1}{2}$.

\textbf{Method 2 — Algebraic trick:}
$\Prob(\text{even heads}) = \dfrac{1+(1-2p)^n}{2}\Big|_{p=1/2} = \dfrac{1}{2}$ for all $n\geq1$.

\textbf{Method 3 — Recurrence:}
$a_n = \tfrac{1}{2}$ satisfies $a_n = p(1-a_{n-1})+(1-p)a_{n-1}$ when $a_{n-1}=\tfrac{1}{2}$.

\textbf{Key takeaway:} Parity of heads is always $\tfrac{1}{2}$ for any $n\geq1$ fair coin tosses.
\end{solutionbox}

\section{Complementary Counting}

\begin{techniquebox}{Complementary Counting}
$\Prob(A)=1-\Prob(A^c)$. Use when ``at least one'' appears and the complement
(``none'' / ``all different'') is easier to count.
\end{techniquebox}

\begin{freeproblem}{3}{\freelabel}{\difficultytitle{medium}}{quant\_interview\_fundamental\_qdsi\_p26}
\textbf{The Birthday Paradox.}
Assume 365 equally likely birthdays. Find the smallest $n$ such that
$\Prob(\text{at least two people share a birthday})>\tfrac{1}{2}$.
\end{freeproblem}

\begin{solutionbox}
$\Prob(\text{all distinct})=\prod_{k=0}^{n-1}(1-k/365)$.
Approximation $\ln(1-x)\approx-x$: threshold at $n(n-1)>730\ln2\approx506$,
so $\mathbf{n=23}$.

\textbf{Key takeaway:} Pairs grow as $\dbinom{n}{2}$, reaching $\sim365$ when $n\approx23$.
General threshold: $n\approx1.18\sqrt{m}$ for pool size $m$.
\end{solutionbox}

\section{Combinatorial Techniques}

\begin{techniquebox}{Non-Adjacent Selection}
Choose $k$ from $\{1,\ldots,n\}$ with no two adjacent:
$\dbinom{n-k+1}{k}$ ways.
(Substitute $y_i=x_i-(i-1)$ to absorb gap constraints.)
\end{techniquebox}

\begin{freeproblem}{4}{\freelabel}{\difficultytitle{medium}}{citadel\_qdsi\_p26}
\textbf{Probability of Choosing Three Non-Consecutive Integers.}
From $\{1,\ldots,10\}$, pick 3 distinct numbers uniformly at random. What is the
probability no two are adjacent?
\end{freeproblem}

\begin{solutionbox}
Total: $\dbinom{10}{3}=120$. Favourable: $\dbinom{8}{3}=56$.
$\Prob = 56/120 = \mathbf{7/15}$.

\textbf{Key takeaway:} Memorise $\dbinom{n-k+1}{k}$ — it turns a constrained adjacency
problem into an unconstrained one.
\end{solutionbox}

\section{Inclusion-Exclusion Principle}

\begin{conceptbox}{Inclusion-Exclusion Principle (IEP)}
\[
  \Prob\!\Bigl(\bigcup_{i=1}^n A_i\Bigr)
  = \sum_i\Prob(A_i) - \sum_{i<j}\Prob(A_i\cap A_j)
  + \sum_{i<j<k}\Prob(A_i\cap A_j\cap A_k) - \cdots
\]
Signal phrase: \textbf{``at least one''}.
\end{conceptbox}

\begin{freeproblem}{5}{\freelabel}{\difficultytitle{medium}}{quant\_interview\_fundamental\_qdsi\_p3}
\textbf{Probability of Consecutive Sixes in Three Dice Rolls.}
Roll a fair six-sided die 3 times. What is the probability of getting at least one
occurrence of the consecutive pattern $66$?
\end{freeproblem}

\begin{solutionbox}
Let $A=\{\text{rolls 1,2 both 6}\}$, $B=\{\text{rolls 2,3 both 6}\}$.
\[
  \Prob(A\cup B)=\frac{1}{36}+\frac{1}{36}-\frac{1}{216}=\frac{11}{216}\approx5.1\%
\]
\textbf{Key takeaway:} Define $A_i=\{\text{pattern starts at }i\}$, apply IEP.
Intersection $A_i\cap A_{i+1}$ is the longer overlapping pattern.
\end{solutionbox}

\section{Circular Arrangements}

\begin{conceptbox}{Circular Permutations}
Fix one person: $(n-1)!$ total. Block of $k$ friends: $(n-k)!\cdot k!$ arrangements.
\[
  \Prob(\text{$k$ friends consecutive}) = \frac{k!\,(n-k)!}{(n-1)!}
\]
\end{conceptbox}

\begin{freeproblem}{6}{\freelabel}{\difficultytitle{medium}}{quant\_interview\_fundamental\_qdsi\_p70}
\textbf{Adjacent Friends Around a Round Table.}
\textbf{Part 1.} 5 people, 2 friends. $\Prob(\text{friends adjacent})$?
\textbf{Part 2.} Generalise to $n$ people, $k$ friends all consecutive.
\end{freeproblem}

\begin{solutionbox}
\textbf{Part 1.} Fix one friend; 4 positions, 2 adjacent: $\Prob=\mathbf{1/2}$.

\textbf{Part 2.}
$\Prob = \dfrac{k!\,(n-k)!}{(n-1)!}$. Verify $n=5, k=2$: $12/24=1/2$. \checkmark

\textbf{Key takeaway:} Always fix one person first. The block trick handles all
consecutive-group problems.
\end{solutionbox}

\section{Advanced Counting and Geometric Probability}

\subsection{Card Dealing}

\begin{premiumproblem}{7}{\premiumlabel}{\difficultytitle{hard}}{quant\_interview\_fundamental\_qdsi\_p17}
\textbf{Probability Each Player Receives an Ace.}
A 52-card deck is dealt to 4 players (13 cards each). What is the probability every
player receives at least one ace?
\end{premiumproblem}

\begin{solutionbox}
\textbf{Step 1: Total number of equally likely deals.}
A deal assigns 13 cards to each of 4 labeled players:
\[
  |\Omega| = \binom{52}{13}\binom{39}{13}\binom{26}{13}\binom{13}{13}
  = \frac{52!}{(13!)^4}
\]

\textbf{Step 2: Each player must receive exactly 1 ace.}
There is no other way to give all 4 players at least one ace (only 4 aces exist).
The number of ways to assign the 4 distinct aces to the 4 labeled players is:
\[
  4!
\]

\textbf{Step 3: Deal the remaining 48 non-ace cards.}
Each player already has 1 ace and needs 12 more cards from the 48 non-aces:
\[
  \binom{48}{12}\binom{36}{12}\binom{24}{12}\binom{12}{12}
  = \frac{48!}{(12!)^4}
\]

\textbf{Step 4: Favorable deals.}
\[
  |\text{favorable}| = 4!\cdot\frac{48!}{(12!)^4}
\]

\textbf{Step 5: Compute the probability.}
\[
  \Prob(\text{each player has an ace})
  = \frac{4!\cdot\dfrac{48!}{(12!)^4}}{\dfrac{52!}{(13!)^4}}
  = 4!\cdot\frac{48!}{52!}\cdot\frac{(13!)^4}{(12!)^4}
\]

\textbf{Step 6: Simplify.}
\[
  \frac{48!}{52!} = \frac{1}{52\cdot51\cdot50\cdot49},
  \qquad
  \left(\frac{13!}{12!}\right)^4 = 13^4
\]
\[
  \Prob(\text{each player has an ace})
  = \frac{4!\cdot 13^4}{52\cdot51\cdot50\cdot49}
\]
Using $52 = 4\cdot13$ to cancel with $4! = 24$:
\[
  = \frac{24\cdot13^4}{4\cdot13\cdot51\cdot50\cdot49}
  = \frac{6\cdot13^3}{51\cdot50\cdot49}
\]

\[
  \boxed{
  \Prob(\text{each player has an ace})
  = \frac{24\cdot 13^4}{52\cdot 51\cdot 50\cdot 49}
  = \frac{6\cdot 13^3}{51\cdot 50\cdot 49}
  = \frac{2197}{20825}
  \approx 10.5\%
  }
\]

\textbf{Key takeaway:} The direct-counting approach --- assign aces first,
then deal non-aces --- is cleaner than IEP here. Despite 4 aces for 4 players,
the probability each receives one is only $\approx 1$ in $9.5$, because most
deals concentrate two or more aces in the same hand.
\end{solutionbox}

\subsection{Derangements}

\begin{keyresult}{Derangement Formula}
$D_n = n!\sum_{k=0}^{n}(-1)^k/k!\approx n!/e$.
For $n\geq5$: $\Prob(\text{derangement})\approx1/e\approx36.8\%$.
\end{keyresult}

\begin{premiumproblem}{8}{\premiumlabel}{\difficultytitle{medium}}{quant\_interview\_fundamental\_qdsi\_p25}
\textbf{The Derangement Problem.}
5 firms, 5 cover letters placed randomly. Probability no letter goes to the correct firm?
\end{premiumproblem}

\begin{solutionbox}
\[
  \Prob = \sum_{k=0}^{5}\frac{(-1)^k}{k!}
  = 1-1+\tfrac{1}{2}-\tfrac{1}{6}+\tfrac{1}{24}-\tfrac{1}{120}
  = \frac{44}{120}=\frac{11}{30}\approx36.7\%
\]
\textbf{Key takeaway:} $\approx1/e$ for $n\geq5$.
Recurrence: $D_n=(n-1)(D_{n-1}+D_{n-2})$.
\end{solutionbox}

\subsection{Geometric Probability}

\begin{techniquebox}{Geometric Probability}
$\Prob(A)=\text{measure of favourable region}/\text{total measure}$.
Identify the right sample space, map the event geometrically, then compute.
\textbf{Always draw the picture.}
\end{techniquebox}

\begin{premiumproblem}{9}{\premiumlabel}{\difficultytitle{hard}}{quant\_interview\_fundamental\_qdsi\_p24}
\textbf{Probability All Random Points Lie in a Semicircle.}
Pick $N$ points uniformly on a circle's circumference. What is the probability
all $N$ lie within some semicircle?
\end{premiumproblem}

\begin{solutionbox}
Let $A_i=\{\text{all other }N-1\text{ points in the semicircle clockwise from }i\}$.
Events are mutually exclusive, $\Prob(A_i)=(1/2)^{N-1}$.
\[
  \boxed{\Prob = \frac{N}{2^{N-1}}}
\]
\textbf{Key takeaway:} Mutual exclusivity of $A_1,\ldots,A_N$ is the key step.
\end{solutionbox}

\begin{premiumproblem}{10}{\premiumlabel}{\difficultytitle{hard}}{quant\_interview\_fundamental\_qdsi\_p10}
\textbf{Probability Two Random Chords Intersect.}
Choose chords by picking four endpoints uniformly on a circle. What is the
probability the chords intersect?
\end{premiumproblem}

\begin{solutionbox}
4 random points $\to$ 3 equally likely pairings. Chords intersect iff endpoints
alternate; exactly 1 of 3 pairings does this.
\[
  \boxed{\Prob = \frac{1}{3}}
\]
\begin{warningbox}
Different chord-sampling methods give different answers (Bertrand's paradox).
Always clarify the sampling method in an interview.
\end{warningbox}
\end{solutionbox}

\begin{premiumproblem}{11}{\premiumlabel}{\difficultytitle{veryhard}}{quant\_interview\_fundamental\_qdsi\_p86}
\textbf{Triangle Formation from Random Stick Breaks.}
\textbf{(1)} Unit stick cut at two uniform points. $\Prob(\text{triangle})$?
\textbf{(2)} First cut uniform; break the longer piece at a uniform point. $\Prob(\text{triangle})$?
\end{premiumproblem}

\begin{solutionbox}
\textbf{(1)} Event = each piece $<1/2$ = two triangles in $[0,1]^2$, total area $1/4$.
$\boxed{\Prob_1=1/4}$.

\textbf{(2)} Condition on first cut $U\sim\mathrm{Unif}(0,\tfrac{1}{2})$, second
$V\sim\mathrm{Unif}(U,1)$. Need $V>1/2$ and $V<U+1/2$:
\[
  \Prob_2 = 2\int_0^{1/2}\frac{U}{1-U}\,dU = 2\ln2-1\approx38.6\%
\]
\textbf{Key takeaway:} Breaking the longer piece raises the probability from
25\% to $\approx$38.6\%.
\end{solutionbox}

\section{Additional Exercises}

\begin{premiumproblem}{A1}{\premiumlabel}{\difficultytitle{hard}}{citadel\_qdsi\_p24}
\textbf{Average Chord Length of the Unit Circle.}
What is the expected distance between two randomly chosen points on a unit circle?

\textit{Hint: chord length $= 2|\sin(\Delta\theta/2)|$; integrate over $\Delta\theta\sim\mathrm{Unif}(0,\pi)$.}
\end{premiumproblem}

\vspace{4pt}

\begin{premiumproblem}{A2}{\premiumlabel}{\difficultytitle{medium}}{quant\_interview\_fundamental\_qdsi\_p69}
\textbf{Balanced Assignment of Seniors Across Teams.}
16 people (4 seniors, 12 juniors) partitioned into 4 labeled groups of 4.
$\Prob(\text{each group has exactly 1 senior})$?

\textit{Hint: count ways to assign 4 seniors to 4 groups, one per group.}
\end{premiumproblem}

\section{Technique Summary}

\begin{tcolorbox}[breakable, colback=EpiNavy!6, colframe=EpiNavy!50,
                  arc=4pt, left=10pt, right=10pt, top=8pt, bottom=8pt]
\textbf{\color{EpiNavy}Topic 1 Technique Toolkit}\vspace{6pt}

\begin{tabular}{@{}p{4.4cm}p{9.6cm}@{}}
\textbf{Uniform probability}    & Count favourable / count total. \\[3pt]
\textbf{Complementary}          & $\Prob(A)=1-\Prob(A^c)$. Use for ``at least one''. \\[3pt]
\textbf{Symmetry}               & Fair coin/dice: use symmetry before computing. \\[3pt]
\textbf{IEP}                    & $\Prob(\bigcup A_i)=\sum\Prob(A_i)-\sum\Prob(A_i\cap A_j)+\cdots$ \\[3pt]
\textbf{Non-adjacent}           & $k$ from $n$, no two adjacent: $\dbinom{n-k+1}{k}$. \\[3pt]
\textbf{Circular}               & Fix one; block of $k$: $(n-k)!\cdot k!$. \\[3pt]
\textbf{Derangement}            & $\Prob\approx1/e\approx36.8\%$ for $n\geq5$. \\[3pt]
\textbf{Geometric}              & Map to region; compute area. Draw it. \\[3pt]
\textbf{Chord intersection}     & 4 pts $\to$ 3 pairings, 1 crosses: $\Prob=1/3$. \\[3pt]
\textbf{Semicircle containment} & $\Prob=N/2^{N-1}$. \\
\end{tabular}
\end{tcolorbox}

\section{Problem Index}

\begin{center}\small
\begin{tabular}{cllll}
\toprule
\textbf{\#} & \textbf{Problem} & \textbf{Tier} & \textbf{Difficulty} & \textbf{Key Technique} \\
\midrule
W1  & Majority Heads (3 coin tosses)     & \freelabel    & \difficulty{easy}     & Symmetry \\
2   & Parity of Heads (10 coins)          & \freelabel    & \difficulty{medium}   & Algebraic trick \\
3   & The Birthday Paradox                & \freelabel    & \difficulty{medium}   & Complementary counting \\
4   & Three Non-Consecutive Integers      & \freelabel    & \difficulty{medium}   & Non-adjacent bijection \\
5   & Consecutive Sixes in Dice Rolls     & \freelabel    & \difficulty{medium}   & IEP \\
6   & Adjacent Friends at Round Table     & \freelabel    & \difficulty{medium}   & Circular permutations \\
\midrule
7   & Each Player Receives an Ace         & \premiumlabel & \difficulty{hard}     & IEP, card counting \\
8   & The Derangement Problem             & \premiumlabel & \difficulty{medium}   & Derangement formula \\
9   & All Points in a Semicircle          & \premiumlabel & \difficulty{hard}     & Mutual exclusivity \\
10  & Two Random Chords Intersect         & \premiumlabel & \difficulty{hard}     & Symmetry, 3 pairings \\
11  & Triangle from Stick Breaks          & \premiumlabel & \difficulty{veryhard} & Geometric probability \\
\midrule
A1  & Average Chord Length                & \premiumlabel & \difficulty{hard}     & Geometric probability \\
A2  & Balanced Senior Assignment          & \premiumlabel & \difficulty{medium}   & Combinatorial counting \\
\bottomrule
\end{tabular}
\end{center}

\epigramfooter

\begin{idtable}
  \idrow{W1}{Majority Heads in Three Coin Tosses}{citadel\_qdsi\_p28}
  \idrow{2}{Parity of Heads in Coin Tosses}{quant\_interview\_fundamental\_qdsi\_p40}
  \idrow{3}{The Birthday Paradox}{quant\_interview\_fundamental\_qdsi\_p26}
  \idrow{4}{Three Non-Consecutive Integers}{citadel\_qdsi\_p26}
  \idrow{5}{Consecutive Sixes in Three Dice Rolls}{quant\_interview\_fundamental\_qdsi\_p3}
  \idrow{6}{Adjacent Friends Around a Round Table}{quant\_interview\_fundamental\_qdsi\_p70}
  \idrow{7}{Each Player Receives an Ace}{quant\_interview\_fundamental\_qdsi\_p17}
  \idrow{8}{The Derangement Problem}{quant\_interview\_fundamental\_qdsi\_p25}
  \idrow{9}{All Points in a Semicircle}{quant\_interview\_fundamental\_qdsi\_p24}
  \idrow{10}{Two Random Chords Intersect}{quant\_interview\_fundamental\_qdsi\_p10}
  \idrow{11}{Triangle Formation from Stick Breaks}{quant\_interview\_fundamental\_qdsi\_p86}
  \idrow{A1}{Average Chord Length of the Unit Circle}{citadel\_qdsi\_p24}
  \idrow{A2}{Balanced Assignment of Seniors Across Teams}{quant\_interview\_fundamental\_qdsi\_p69}
\end{idtable}

\end{document}